\section{Step 79 Check boundaries before comparing work}\label{h263-writer-s09}
Run public feedback on your current source:
\begin{verbatim}
python3 construction/h263-bit-writer-construction/grade.py \
    my-writer.asm --cli /path/to/mars-cli.jar \
    --report writer-results.json
\end{verbatim}
The checker constructs an independent bit sequence and compares the affected
bytes. It checks both simulator byte orders, every starting bit offset, representative
widths including zero and 32, exact fits, insufficient capacity, wide malformed
controls, guarded context and the maximum-capacity endpoint. Read/write traces
reject access outside bytes touched by the field.

Measure entry-to-return instruction work for the same value, width and initial
cursor in both implementations. Include validation, loop branches, byte accesses
and the committed count update. This local work comparison is not a whole-encoder
wall-time measurement or a grading threshold.

For the named all-one 32-bit field starting at bit three, the tested bit-at-a-time
entry executes 656 instructions and the byte-fragment entry executes 189, including
admission and the final counter store. Both simulator byte orders give the same
counts. This saves 467 instructions for that append; it does not measure an entire
coefficient event, file write or encoder launch. Different valid implementations
need not reproduce these counts.

\paragraph{Observe} A bit-at-a-time implementation remains valid after the
byte-fragment implementation is introduced. Correct output does not require a
particular instruction count or private label.
\Needspace{8\baselineskip}
\paragraph{Check} What must match when comparing the two schedules?
\begin{itemize}
\item[A.] Only their scratch-register names.
\item[B.] The host computer's display theme.
\item[C.] Field value, width, initial bit cursor and complete measured boundary.
\end{itemize}
